\section{AMC 8 1995}
        \begin{enumerate}
        	\item (\textit{AMC 8 1995}) Walter has exactly one penny, one nickel, one dime and one quarter in his pocket.  What percent of one dollar is in his pocket?


 $\text{(A)}\ 4\% \qquad \text{(B)}\ 25\% \qquad \text{(C)}\ 40\% \qquad \text{(D)}\ 41\% \qquad \text{(E)}\ 59\%$

 \textbf{Answer: }D



\textbf{Solution}

Walter has $1 + 5 + 10 + 25 = 41$ cents in his pocket.  There are $100$ cents in a dollar.  Therefore, Walter has $\frac{41}{100}\cdot 100\% = 41\%$ of a dollar in his pocket, and the right answer is $\boxed{D}$.

 


	\item (\textit{AMC 8 1995}) Jose is $4$ years younger than Zack.  Zack is $3$ years older than Inez.  Inez is $15$ years old.  How old is Jose?


 $\text{(A)}\ 8 \qquad \text{(B)}\ 11 \qquad \text{(C)}\ 14 \qquad \text{(D)}\ 16 \qquad \text{(E)}\ 22$

 \textbf{Answer: }C



\textbf{Solution}

Zack is $3$ years older than Inez, who is $15$.  Therefore, Zack is $15 + 3 = 18$ years old.

 Jose is $4$ years younger than Zack, who is $18$.  Therefore, Jose is $18 - 4 = 14$ years old, and the answer is $\boxed{C}$.

 


	\item (\textit{AMC 8 1995}) Which of the following operations has the same effect on a number as multiplying by $\dfrac{3}{4}$ and then dividing by $\dfrac{3}{5}$?


 $\text{(A)}\ \text{dividing by }\dfrac{4}{3} \qquad \text{(B)}\ \text{dividing by }\dfrac{9}{20} \qquad \text{(C)}\ \text{multiplying by }\dfrac{9}{20} \qquad \text{(D)}\ \text{dividing by }\dfrac{5}{4} \qquad \text{(E)}\ \text{multiplying by }\dfrac{5}{4}$

 \textbf{Answer: }E



\textbf{Solution}

Dividing by $\frac{3}{5}$ is the same as multiplying by its reciprocal, $\frac{5}{3}$.

 Multiplying a number by two numbers is the same as multiplying the number by the product of those two numbers.  Thus, multiplying by $\frac{3}{4}$ and then $\frac{5}{3}$ is the same as multiplying by $\frac{3\cdot 5}{4\cdot 3} = \frac{5}{4}$, and the answer is $\boxed{E}$.

 


	\item (\textit{AMC 8 1995}) A teacher tells the class,


 
\textit{"Think of a number, add 1 to it, and double the result.  Give the answer to your partner.  Partner, subtract 1 from the number you are given and double the result to get your answer."}


 


 Ben thinks of $6$, and gives his answer to Sue.  What should Sue's answer be?


 $\text{(A)}\ 18 \qquad \text{(B)}\ 24 \qquad \text{(C)}\ 26 \qquad \text{(D)}\ 27 \qquad \text{(E)}\ 30$

 \textbf{Answer: }C



\textbf{Solution}

Ben adds $1$ to $6$, getting $1 + 6 = 7$.  Then Ben doubles it, to get $7\cdot 2 = 14$.  He gives this number to Sue.  Then Sue subtracts $1$ from the number, to get $14 - 1 = 13$.  Finally, Sue doubles the result to get $13\cdot 2 = 26$, and the answer is $\boxed{C}$.

 


	\item (\textit{AMC 8 1995}) Find the smallest whole number that is larger than the sum


 
 \[2\dfrac{1}{2}+3\dfrac{1}{3}+4\dfrac{1}{4}+5\dfrac{1}{5}.\]
$\text{(A)}\ 14 \qquad \text{(B)}\ 15 \qquad \text{(C)}\ 16 \qquad \text{(D)}\ 17 \qquad \text{(E)}\ 18$

 \textbf{Answer: }C



\textbf{Solution 1}

Adding the whole numbers gives $2 + 3 + 4 + 5 = 14$.

 Adding the fractions gives $\frac{1}{2} + \frac{1}{3} + \frac{1}{4} + \frac{1}{5} = \frac{30}{60} + \frac{20}{60} + \frac{15}{60} + \frac{12}{60} = \frac{77}{60}$.  This will create one more whole, and a fraction that is less than $1$.  Thus, the smallest whole number that is less than $15$ plus some fractional part is $16$, and the answer is $\boxed{C}$.

 



\textbf{Solution 2}

Convert the fractional parts to decimals, and approximate the answer.  $2.5 + 3.33 + 4.25 + 5.2 = 15.28$, and the answer is $16$, which is $\boxed{C}$.

 



\textbf{Solution 3 (Speedy)}

Clearly, break the mixed numbers into $2+3+4+5+\frac{1}{2} + \frac{1}{3} + \frac{1}{4} + \frac{1}{5}$. 

 This is the same as $14+(\frac{1}{2} + \frac{1}{3}) + (\frac{1}{4} + \frac{1}{5})$. Notice that clearly, $\frac{1}{2}+\frac{1}{3}=\frac{5}{6}$, and $1>\frac{1}{4}+\frac{1}{5}>\frac{1}{6}$ (clearly, $\frac{1}{4}>\frac{1}{6}$). This means that the fractional parts is between 1 and 2, so when that is added to 14, the original expression is between 15, and 16. So the least integer that is bigger than this sum is clearly 16. So select $\boxed{C}$.

 


	\item (\textit{AMC 8 1995}) Figures $I$, $II$, and $III$ are squares.  The perimeter of $I$ is $12$ and the perimeter of $II$ is $24$.  The perimeter of $III$ is


 
                    \begin{figure}[H]
                        \centering
                        \includegraphics[width=0.5\textwidth]{Problem/AMC_8_AJHSME_1995_6_Problem_Figure_1.png}
                    \end{figure}
                    
$\text{(A)}\ 9 \qquad \text{(B)}\ 18 \qquad \text{(C)}\ 36 \qquad \text{(D)}\ 72 \qquad \text{(E)}\ 81$

 \textbf{Answer: }C



\textbf{Solution 1}

Since the perimeter of $I$ $12$, each side is $\frac{12}{4} = 3$.  

 Since the perimeter of $II$ is $24$, each side is $\frac{24}{4} = 6$.

 The side of $III$ is equal to the sum of the sides of $I$ and $II$.  Therefore, the side of $III$ is $3 + 6 = 9$.

 Since $III$ is also a square, it has an perimeter of $9\cdot 4 = 36$, and the answer is $\boxed{C}$.

 



\textbf{Solution 2}

Let a side of $I$ equal $x$, and let a side of $II$ equal $y$. The perimeter of $I$ is $4x$, and the perimeter of $II$ is $4y$. One side of $III$ has length $x+y$, so the perimeter is $4x+4y$, which just so happens to be the sum of the perimeters of $I$ and $II$, giving us $12+24=36$, or answer $\boxed{C}$.

 


	\item (\textit{AMC 8 1995}) At Clover View Junior High, one half of the students go home on the school bus.  One fourth go home by automobile.  One tenth go home on their bicycles.  The rest walk home.  What fractional part of the students walk home?


 $\text{(A)}\ \dfrac{1}{16} \qquad \text{(B)}\ \dfrac{3}{20} \qquad \text{(C)}\ \dfrac{1}{3} \qquad \text{(D)}\ \dfrac{17}{20} \qquad \text{(E)}\ \dfrac{9}{10}$

 \textbf{Answer: }B



\textbf{Solution}

The students who walk home are

 $1 - \frac{1}{2} - \frac{1}{4} - \frac{1}{10}$
$\frac{1}{2} - \frac{1}{4} - \frac{1}{10}$
$\frac{1}{4} - \frac{1}{10}$
$\frac{5}{20} - \frac{2}{20}$
$\frac{3}{20}$ of the entire student population.  Therefore, the answer is $\boxed{B}$. 

 


	\item (\textit{AMC 8 1995}) An American traveling in Italy wishes to exchange American money (dollars) for Italian money (lire).  If 3000 lire = $1.60, how much lire will the traveler receive in exchange for $1.00?


 $\text{(A)}\ 180 \qquad \text{(B)}\ 480 \qquad \text{(C)}\ 1800 \qquad \text{(D)}\ 1875 \qquad \text{(E)}\ 4875$

 \textbf{Answer: }D



\textbf{Solution}

$\frac{3000 \text{ lire}}{1.6 \text{ dollars}} = \frac{x \text{ lire}}{1 \text{ dollar}}$
$x = 1875$, and the answer is $\boxed{D}$



	\item (\textit{AMC 8 1995}) Three congruent circles with centers $P$, $Q$, and $R$ are tangent to the sides of rectangle $ABCD$ as shown.  The circle centered at $Q$ has diameter $4$ and passes through points $P$ and $R$.  The area of the rectangle is


 
                    \begin{figure}[H]
                        \centering
                        \includegraphics[width=0.5\textwidth]{Problem/AMC_8_AJHSME_1995_9_Problem_Figure_1.png}
                    \end{figure}
                    
$\text{(A)}\ 16 \qquad \text{(B)}\ 24 \qquad \text{(C)}\ 32 \qquad \text{(D)}\ 64 \qquad \text{(E)}\ 128$

 \textbf{Answer: }C



\textbf{Solution 1}

If circle $Q$ has diameter $4$, then so do congruent circles $P$ and $R$.  Draw a diameter through $P$ parallel to $AD$.  The diameter will be congruent to $AD$, and thus $AD = 4$, which is the height of the rectangle.

 Draw a horizontal line $PQR$ that extends to the sides of the rectangle.  This line is $2$ diameters long, so it has length $4\cdot 2 = 8$.  It is parallel and congruent to $AB$, so the width of the rectangle is $8$.

 Thus, the area of the rectangle is $lw = 4\cdot 8 = 32$, and the answer is $\boxed{C}$.

 



\textbf{Solution 2}

The area of circle $P$ and circle $R$ almost fill the rectangle.  Circle $P$ has radius $\frac{4}{2} = 2$, and so does circle $R$.  Thus, the sum of their areas is $\pi \cdot 2^2 + \pi \cdot 2^2 = 8\pi$.  Since $\pi \approx 3.14$, the area of the circles is just over $24$.  The area of the rectangle is greater than the area of the circles.  $64$ would be too high, as the two circles appear to take up much more than half the area of the rectangle.  But $24$ is too low, as it is less than the area of the two circles.  Thus, the only reasonable answer is $\boxed{C}$, which is $32$.

 


	\item (\textit{AMC 8 1995}) A jacket and a shirt originally sold for 80 dollars and 40 dollars, respectively.  During a sale Chris bought the 80 dollar jacket at a $40\%$ discount and the 40 dollar shirt at a $55\%$ discount.  The total amount saved was what percent of the total of the original prices?


 $\text{(A)}\ 45\% \qquad \text{(B)}\ 47\dfrac{1}{2}\% \qquad \text{(C)}\ 50\% \qquad \text{(D)}\ 79\dfrac{1}{6}\% \qquad \text{(E)}\ 95\%$.


 
 \textbf{Answer: }A



\textbf{Solution}

Chris would have spent $40 + 80 = 120$ dollars if there were no sale.

 The sale means that he saved $80(0.40) + 40(0.55) = 32 + 22 = 54$ dollars.

 He saved $\frac{54}{120}\cdot 100\% = 0.45 \cdot 100\% = 45\%$ off the total of the original prices, and the answer is $\boxed{A}$.

 


	\item (\textit{AMC 8 1995}) Jane can walk any distance in half the time it takes Hector to walk the same distance.  They set off in opposite directions around the outside of the 18-block area as shown.  When they meet for the first time, they will be closest to


 
                    \begin{figure}[H]
                        \centering
                        \includegraphics[width=0.5\textwidth]{Problem/AMC_8_AJHSME_1995_11_Problem_Figure_1.png}
                    \end{figure}
                    
$\text{(A)}\ A \qquad \text{(B)}\ B \qquad \text{(C)}\ C \qquad \text{(D)}\ D \qquad \text{(E)}\ E$

 \textbf{Answer: }D



\textbf{Solution 1}

Counting around, when Jane walks $12$ steps, she will be at $D$.  When Hector walks $6$ steps, he will also be at $D$.  Since Jane has walked twice as many steps as Hector, they will reach this spot at the same time.  Thus, the answer is $\boxed{D}$.

 



\textbf{Solution 2}

We count to find that the perimeter of the rectangle is 18. We divide by 3 to see how far Hector has walked when he meets Jane, this distance is 6. We then move 6 steps counterclockwise to find the point they met or $\boxed{D}$.

 


	\item (\textit{AMC 8 1995}) A \textit{lucky} year is one in which at least one date, when written in the form month/day/year, has the following property: \textit{The product of the month times the day equals the last two digits of the year}.  For example, 1956 is a lucky year because it has the date 7/8/56 and $7\times 8 = 56$.  Which of the following is NOT a lucky year?


 $\text{(A)}\ 1990 \qquad \text{(B)}\ 1991 \qquad \text{(C)}\ 1992 \qquad \text{(D)}\ 1993 \qquad \text{(E)}\ 1994$

 \textbf{Answer: }E



\textbf{Solution}

We examine only the factors of $90, 91, 92, 93,$ and $94$ that are less than $13$, because for a year to be lucky, it must have at least one factor between $1$ and $12$ to represent the month.

 $90$ has factors of $1, 2, 3, 5, 6, 9,$ and $10$.  Dividing $90$ by the last number $10$, gives $9$.  Thus, $10/9/90$ is a valid date, and $1990$ is a lucky year.

 $91$ has factors of $1$ and $7$.  Dividing $91$ by $7$ gives $13$, and $7/13/91$ is a valid date.  Thus, $1991$ is a lucky year.

 $92$ has factors of $1, 2$ and $4$.  Dividing $92$ by $4$ gives $23$, and $4/23/92$ is a valid date.  Thus, $1992$ is a lucky year.

 $93$ has factors of $1$ and $3$.  Dividing $93$ by $3$ gives $31$, and $3/31/93$ is a valid date.  Thus, $1993$ is a lucky year.

 $94$ has factors of $1$ and $2$.  Dividing $94$ by $2$ gives $47$, and no month has $47$ days.  Thus, $1994$ is not a lucky year, and the answer is $\boxed{E}$.

 


	\item (\textit{AMC 8 1995}) In the figure, $\angle A$, $\angle B$, and $\angle C$ are right angles.  If $\angle AEB = 40^\circ$ and $\angle BED = \angle BDE$, then $\angle CDE =$

                    \begin{figure}[H]
                        \centering
                        \includegraphics[width=0.5\textwidth]{Problem/AMC_8_AJHSME_1995_13_Problem_Figure_1.png}
                    \end{figure}
                    
$\text{(A)}\ 75^\circ \qquad \text{(B)}\ 80^\circ \qquad \text{(C)}\ 85^\circ \qquad \text{(D)}\ 90^\circ \qquad \text{(E)}\ 95^\circ$

 \textbf{Answer: }E



\textbf{Solution}

Because $\angle BED=\angle BDE$, $\angle B=90^\circ$, and $\triangle BED$ is a triangle, we get: 
 \[\angle B + \angle BED +\angle BDE=180\] 
 \[90 +\angle BED +\angle BED=180\] 
 \[2\angle BED=90\] 
 \[\angle BED=\angle BDE=45^\circ\]

 So $\angle AED=\angle AEB +\angle BED=40 +45=85^\circ$ Since ACDE is a quadrilateral, the sum of its angles is 360. Therefore: 
 \[\angle A +\angle C +\angle CDE +\angle AED=360\] 
 \[90 +90 +\angle CDE +85=360\] 
 \[\angle CDE=95^\circ \text{(E)}\]



	\item (\textit{AMC 8 1995}) A team won 40 of its first 50 games.  How many of the remaining 40 games must this team win so it will have won exactly 70\% of its games for the season?


 $\text{(A)}\ 20 \qquad \text{(B)}\ 23 \qquad \text{(C)}\ 28 \qquad \text{(D)}\ 30 \qquad \text{(E)}\ 35$

 \textbf{Answer: }B



\textbf{Solution 1}

The team will play a total of $50+40=90$ games, so to finish with a $70\%$ win percentage it must win $0.7 \cdot 90=63$ games. Therefore, it must win $63-40=\boxed{\textbf{(B)}\ 23}$ of the remaining games.

 



\textbf{Solution 2}

Let $x$ be the number of wins in the team's remaining $40$ games.  Then the team wins $40+x$ of its $90$ games.

 Since $70\% = \frac{7}{10}$, it follows that $\frac{40+x}{90}=\frac{7}{10}$.  Multiplying both sides of this equation by $90$ yields $40 + x = 63$, so $x=\boxed{\textbf{(B)}\ 23}$.

 


	\item (\textit{AMC 8 1995}) What is the $100^\text{th}$ digit to the right of the decimal point in the decimal form of $4/37$?


 $\text{(A)}\ 0 \qquad \text{(B)}\ 1 \qquad \text{(C)}\ 2 \qquad \text{(D)}\ 7 \qquad \text{(E)}\ 8$

 \textbf{Answer: }B



\textbf{Solution}

$\frac{4}{37}=\frac{12}{111}=\frac{108}{999}=0.108108108...$

 Since this repeats every three digits, digit number x = digit number (x mod 3), and the 100th digit = (100 mod 3)th digit = 1st digit = $\boxed{\text{(B)}\ 1}$



	\item (\textit{AMC 8 1995}) Students from three middle schools worked on a summer project.


 
            \begin{itemize}
        \item Seven students from Allen school worked for 3 days.

        \item Four students from Balboa school worked for 5 days.

        \item Five students from Carver school worked for 9 days.
            \end{itemize}
            The total amount paid for the students' work was $774$.  Assuming each student received the same amount for a day's work, how much did the students from Balboa school earn altogether?


 $\text{(A)}\ 9.00\text{ dollars} \qquad \text{(B)}\ 48.38\text{ dollars} \qquad \text{(C)}\ 180.00\text{ dollars} \qquad \text{(D)}\ 193.50\text{ dollars} \qquad \text{(E)}\ 258.00\text{ dollars}$

 \textbf{Answer: }C



\textbf{Solution}

The students combined for $7\cdot 3 + 4 \cdot 5 + 5 \cdot 9 = 86$ student workdays.  Therefore, each student earned $\$774 \div 86 = \$9$ per workday.  Since the Balboa students completed $5 \cdot 4 = 20$ workdays, they earned a total of $20 \cdot \$9 = \boxed{\textbf{(C)}\ \$180.00}$. 

 


	\item (\textit{AMC 8 1995}) The table below gives the percent of students in each grade at Annville and Cleona elementary schools:


 
 \[\begin{tabular}{rccccccc}  & \textbf{\underline{K}} & \textbf{\underline{1}} & \textbf{\underline{2}} & \textbf{\underline{3}} & \textbf{\underline{4}} & \textbf{\underline{5}} & \textbf{\underline{6}} \\ \textbf{Annville:} & 16\% & 15\% & 15\% & 14\% & 13\% & 16\% & 11\% \\ \textbf{Cleona:} & 12\% & 15\% & 14\% & 13\% & 15\% & 14\% & 17\%  \end{tabular}\]


 Annville has 100 students and Cleona has 200 students.  In the two schools combined, what percent of the students are in grade 6?


 $\text{(A)}\ 12\% \qquad \text{(B)}\ 13\% \qquad \text{(C)}\ 14\% \qquad \text{(D)}\ 15\% \qquad \text{(E)}\ 28\%$

 \textbf{Answer: }D



\textbf{Solution}

By the tables, Annville has $11$ 6th graders and Cleona has $34$. Together they have $45$ 6th graders and $300$ total students, so the percent is $\frac{45}{300}=\frac{15}{100}= \boxed{\text{(D)}\ 15\%}$



	\item (\textit{AMC 8 1995}) The area of each of the four congruent L-shaped regions of this 100-inch by 100-inch square is 3/16 of the total area.  How many inches long is the side of the center square?


 
                    \begin{figure}[H]
                        \centering
                        \includegraphics[width=0.5\textwidth]{Problem/AMC_8_AJHSME_1995_18_Problem_Figure_1.png}
                    \end{figure}
                    
$\text{(A)}\ 25 \qquad \text{(B)}\ 44 \qquad \text{(C)}\ 50 \qquad \text{(D)}\ 62 \qquad \text{(E)}\ 75$

 \textbf{Answer: }C



\textbf{Solution}

The area taken up by the L's is $4*\frac{3}{16}=\frac{3}{4}$ of the area of the whole square. What remains has $\frac{1}{4}$ of the area of the larger square. $\frac{100*100}{4}=\frac{100}{2}*\frac{100}{2}=50*50$ is the area of the smaller square, so its side length is 50. $\text{(C)}$



	\item (\textit{AMC 8 1995}) The graph shows the distribution of the number of children in the families of the students in Ms. Jordan's English class.  The median number of children in the family for this distribution is


 
                    \begin{figure}[H]
                        \centering
                        \includegraphics[width=0.5\textwidth]{Problem/AMC_8_AJHSME_1995_19_Problem_Figure_1.png}
                    \end{figure}
                    
$\text{(A)}\ 1 \qquad \text{(B)}\ 2 \qquad \text{(C)}\ 3 \qquad \text{(D)}\ 4 \qquad \text{(E)}\ 5$

 \textbf{Answer: }D



\textbf{Solution}

Counting, there are thirteen total families. The middle number is 7th in either direction, and it is easy to see from the right side that this number is 4 $\text{(D)}$



	\item (\textit{AMC 8 1995}) Diana and Apollo each roll a standard die obtaining a number at random from 1 to 6.  What is the probability that Diana's number is larger than Apollo's number?


 $\text{(A)}\ \dfrac{1}{3} \qquad \text{(B)}\ \dfrac{5}{12} \qquad \text{(C)}\ \dfrac{4}{9} \qquad \text{(D)}\ \dfrac{17}{36} \qquad \text{(E)}\ \dfrac{1}{2}$

 \textbf{Answer: }B



\textbf{Solution 1}

Note that the probability of Diana rolling a number larger than Apollo's is the same as the probability of Apollo's being more than Diana's. If we denote this common probability $D$, then $2D+P($Apollo=Diana$)=1$. Now all we need to do is find $P($Apollo=Diana$)$. There are $6(6)=36$ possibilities total, and 6 of those have Apollo=Diana, so $P($Apollo=Diana$)=\frac{6}{36}=\frac{1}{6}$. Going back to our first equation and solving for D, we get 
 \[2D+\frac{1}{6}=1\] 
 \[2D=\frac{5}{6}\] 
 \[D=\frac{5}{12} \Rightarrow \mathrm{(B)}\]




\textbf{Solution 2}

We can use simple casework to solve this problem too. There are six cases based on Apollo's Roll.Apollo Rolls a 1: Diana could roll a $2$, $3$, $4$, $5$, or $6$.Apollo Rolls a 2: Diana could roll a $3$, $4$, $5$, or $6$.Apollo Rolls a 3: Diana could roll a $4$, $5$, or $6$.Apollo Rolls a 4: Diana could roll a $5$ or $6$.Apollo Rolls a 5: Diana could roll a $6$.Apollo Rolls a 6: There are no successful outcomes.The total amount of successful cases is $5+4+3+2+1 = 15$. The total amount of possible cases is $6(6) = 36$. Therefore, the probability of Diana rolling a bigger number is $\frac{15}{36} = \frac{5}{12} \Rightarrow \mathrm{(B)}$




\textbf{Solution 3 (2 seconds, quick)}

Define a "win" for Player 1 as rolling a greater number than Player 2. Because potential ties make up $\frac{6}{36}=\frac{1}{6}$ of all $36$ combinations, Apollo and Diana "split" the wins in the remaining $1-\frac{1}{6}=\frac{5}{6}$. By symmetry, each person wins half of these, so the probability Diana wins is $\frac{\frac{5}{6}}{2}=\frac{5}{12}$ $\Rightarrow \mathrm{(B)}$



	\item (\textit{AMC 8 1995}) A plastic snap-together cube has a protruding snap on one side and receptacle holes on the other five sides as shown.  What is the smallest number of these cubes that can be snapped together so that only receptacle holes are showing?


 
                    \begin{figure}[H]
                        \centering
                        \includegraphics[width=0.5\textwidth]{Problem/AMC_8_AJHSME_1995_21_Problem_Figure_1.png}
                    \end{figure}
                    
$\text{(A)}\ 3 \qquad \text{(B)}\ 4 \qquad \text{(C)}\ 5 \qquad \text{(D)}\ 6 \qquad \text{(E)}\ 8$

 \textbf{Answer: }B



\textbf{Solution}

It is rather easy to see that the only way that you can possibly have three cubes without protruding snaps is to have them in a triangle formation. However, in this case, the exterior angles are $120^\circ$, whereas the outside angle of the cubes are $90^\circ$ each. A square like formation will work in the case of 4 cubes, so the answer is $\text{(B)}$.

 


	\item (\textit{AMC 8 1995}) The number 6545 can be written as a product of a pair of positive two-digit numbers.  What is the sum of this pair of numbers?


 $\text{(A)}\ 162 \qquad \text{(B)}\ 172 \qquad \text{(C)}\ 173 \qquad \text{(D)}\ 174 \qquad \text{(E)}\ 222$

 \textbf{Answer: }A



\textbf{Solution 1}

The prime factorization of $6545$ is $5\cdot7\cdot11\cdot17  =385\cdot17$, which contains a three digit number, but we want 6545 to be expressed as ab x cd. Now we do trial and error: 
 \[5\cdot7=35 \text{,   } 11\cdot17=187 \text{  X}\] 
 \[5\cdot11=55 \text{,   } 7\cdot17=119 \text{  X}\] 
 \[5\cdot17=85 \text{,   } 7\cdot11=77 \text{ }\surd\] 
 \[85+77= \boxed{\text{(A)}\ 162}\]




\textbf{Solution 2}

Notice that $6545 = 6561 - 16 = 81^2 - 4^2 = (81 - 4)(81 + 4)$. Thus, the sum is $(81 - 4) + (81 + 4) = 2 \times 81 = \boxed{\text{(A)}\ 162}$. 

 This solution is motivated by noticing that a large perfect square, 6561, is close to 6545. In fact, this number is $3^8$, which makes it easily recognized.

 - GaloisTorrent

 



\textbf{Solution 3}

Let $x$ and $y$ be the numbers, with $x \leq y$.  Since $xy = 6545$ is odd, both $x$ and $y$ are odd.

 Since $y < 100$, $x > \frac{6545}{100} = 65.45$.  Thus $66 \leq x, y \leq 99$.

 Since $6545$ is divisible by $5$, one of $x, y$ is divisible by $5$.  Therefore, either $x$ or $y$ is in the set $\{75, 85, 95\}$.  The only divisor of $6545$ in this set is $85$, and $6545 \div 85 = 77$.

 Therefore, $x = 77$ and $y = 85$, so $x + y = \boxed{\textbf{(A)}\ 162}$.

 -j314andrews

 


	\item (\textit{AMC 8 1995}) How many four-digit whole numbers are there such that the leftmost digit is odd, the second digit is even, and all four digits are different?


 $\text{(A)}\ 1120 \qquad \text{(B)}\ 1400 \qquad \text{(C)}\ 1800 \qquad \text{(D)}\ 2025 \qquad \text{(E)}\ 2500$

 \textbf{Answer: }B



\textbf{Solution}

Count from left to right: There are 5 choices for the first digit, 5 choices for the second, 8 remaining choices for the third, and 7 remaining for the fourth, so there are $5*5*8*7= \boxed{\text{(B)}\ 1400}$ numbers.

 


	\item (\textit{AMC 8 1995}) In parallelogram $ABCD$, $\overline{DE}$ is the altitude to the base $\overline{AB}$ and $\overline{DF}$ is the altitude to the base $\overline{BC}$. [\textbf{Note:} \textit{Both pictures represent the same parallelogram.}]  If $DC=12$, $EB=4$, and $DE=6$, then $DF=$

                    \begin{figure}[H]
                        \centering
                        \includegraphics[width=0.5\textwidth]{Problem/AMC_8_AJHSME_1995_24_Problem_Figure_1.png}
                    \end{figure}
                    
$\text{(A)}\ 6.4 \qquad \text{(B)}\ 7 \qquad \text{(C)}\ 7.2 \qquad \text{(D)}\ 8 \qquad \text{(E)}\ 10$

 \textbf{Answer: }C



\textbf{Solution 1}

The area of $ABCD$ is $DC \cdot DE = 12 \cdot 6 = 72$.  Since $DE = 6$ and $AE = AB - BE = 12-4=8$, by the Pythagorean Theorem, $AD = \sqrt{6^2+8^2}=10$.  

 Since the area of $ABCD$ can also be written as $AD \cdot DF$, it follows that $10 \cdot DF = 72$, so $DF = \boxed{\textbf{(C)}\ 7.2}$.

 



\textbf{Solution 2}

Since $AE = 12 - 4 = 8$, by the Pythagorean theorem, $AD = \sqrt{6^2 + 8^2} = 10$.

 Since $\angle A = \angle C$ and $\angle DEA = \angle DFC = 90^{\circ}$, $\triangle DEA \sim \triangle DFC$.  

 Thus $\frac{DE}{DA} = \frac{DF}{DC}$, so $\frac{6}{10} = \frac{DF}{12}$ and thus $DF = \boxed{\textbf{(C)}\ 7.2}$

 -j314andrews

 


	\item (\textit{AMC 8 1995}) Buses from Dallas to Houston leave every hour on the hour.  Buses from Houston to Dallas leave every hour on the half hour.  The trip from one city to the other takes 5 hours.  Assuming the buses travel on the same highway, how many Dallas-bound buses does a Houston-bound bus pass in the highway (not in the station)?


 $\text{(A)}\ 5 \qquad \text{(B)}\ 6 \qquad \text{(C)}\ 9 \qquad \text{(D)}\ 10 \qquad \text{(E)}\ 11$

 \textbf{Answer: }D



\textbf{Solution 1}

Let $B$ be any Houston-bound bus.  Let $X$ be the first bus to arrive in Dallas after $B$ leaves Dallas.  At the time $B$ leaves, $X$ is $0.5$ hours from Dallas and thus left Houston $4.5$ hours before $B$ left Dallas.  Let $Y$ be the last bus to leave Houston before $B$ arrives in Houston.  Then $Y$ must leave Houston $4.5$ hours after $B$ left Dallas.  Therefore, $Y$ leaves Houston $9$ hours after $X$.  Since $B$ must pass every bus between $X$ and $Y$ inclusive, $B$ passes a total of $9 + 1 = \boxed{\text{(D)} \ 10}$ Dallas-bound buses.

 



\textbf{Solution 2}

Say you are on the Houston-bound bus that left at 12:30 in the afternoon, looking out the window to see how many buses you pass. At 12:45 pm, the Dallas bus that left at 8:00 am is 4:45 away (Note - $a:b$ - $a$ is for hrs. and $b$ is for min.) from Dallas, and therefore 15 minutes from Houston. Your bus is also 15 minutes from Houston, so you are delighted at the first bus you have passed. 

 At 1:15 pm, the 9:00 am Dallas bus meets you again, being 4:15 away from Dallas and therefore 45 minutes from Houston. After a while you might notice that the buses meet you in 30 minute intervals after 12:45 pm; indeed, the bus that left an hour later than the bus most recently has traveled 30 minutes closer, and so have you, for a total of an hour closer, and when you passed the most recent bus, the bus after it had indeed left an hour later. You decide to predict how many buses you pass by finding all valid times that are 30 minute intervals from 12:45. They are12:45, 1:15, 1:45, 2:15, 2:45, 3:15, 3:45, 4:15, 4:45, 5:15, and 5:30, your arrival time.Hence, you will pass $\boxed{\text{(D)}\ 10}$ Dallas-bound buses.

 


\end{enumerate}